% J42 — A Discrete sinc^2 Identity in Finite-Dimensional Quantum Mechanics
%
% Brayden R. Sanders, M. Gish, 7Site LLC
% Target venue: Letters in Mathematical Physics (preferred; JMP per-venue cap reached --
%               see cover letter and README §5/§6)
% Source: Theorem 3.1 closed form descends from first_g_sinc2_FINAL.tex (J03 / J04
%         corpus); QM-on-cyclic-group framing is novel for J42.
% Verification: manuscript/verify_J42_sinc2.py  (max deviation 3.33e-16; deterministic;
%               runs in <2 s)
% MSC 2020: 81S05, 42A16, 11A41, 11Y05

\documentclass[11pt,reqno]{amsart}

\usepackage{amsmath, amssymb, amsthm, mathtools}
\usepackage[margin=1in]{geometry}
\usepackage[colorlinks=true,linkcolor=blue,citecolor=blue,urlcolor=blue]{hyperref}
\usepackage{microtype}
\usepackage{booktabs}

\theoremstyle{plain}
\newtheorem{theorem}{Theorem}[section]
\newtheorem{proposition}[theorem]{Proposition}
\newtheorem{lemma}[theorem]{Lemma}
\newtheorem{corollary}[theorem]{Corollary}

\theoremstyle{definition}
\newtheorem{definition}[theorem]{Definition}

\theoremstyle{remark}
\newtheorem*{remark}{Remark}

\newcommand{\Z}{\mathbb{Z}}
\newcommand{\Q}{\mathbb{Q}}
\newcommand{\C}{\mathbb{C}}
\newcommand{\R}{\mathbb{R}}
\newcommand{\spf}{\operatorname{spf}}
\DeclareMathOperator{\sinc}{sinc}

\title[A discrete $\sinc^{2}$ identity in finite-dimensional QM]
{A Discrete $\sinc^{2}$ Identity in Finite-Dimensional Quantum Mechanics}

\author{Brayden R.\ Sanders}
\address{7Site LLC, Hot Springs, Arkansas, USA}
\email{brayden@7site.co}

\author{M.\ Gish}
\address{Independent Researcher, USA}
\email{monica.gish1992@gmail.com}

\subjclass[2020]{81S05, 42A16, 11A41, 11Y05}
\keywords{Fej\'er kernel, finite-dimensional quantum mechanics, cyclic group,
discrete Fourier transform, $\sinc^{2}$, first-zero theorem, coprimality partition}

\date{\today}

\begin{document}

\begin{abstract}
For an integer $f \ge 2$ and $k \ge 1$, define the unit-amplitude exponential sum on the
cyclic position alphabet $\{1,\dots,k\}$ at momentum quantum $1/f$,
\[
  S(k,f) \;=\; \tfrac{1}{k}\sum_{j=1}^{k} e^{2\pi i j / f},
  \qquad R(k,f) \;=\; |S(k,f)|^{2}.
\]
We prove the closed form $R(k,f) = \sin^{2}(\pi k/f)/(k^{2}\sin^{2}(\pi/f))$, identify it
as the Fej\'er kernel restricted to integer arguments, and read it as the squared overlap
between a momentum eigenstate and a normalized position-space rectangular window in finite
QM on $\Z/N\Z$. Three QM-relevant consequences follow: (i) the squared overlap of a
momentum eigenstate with the normalized rectangular position window
(Proposition~\ref{prop:overlap}), (ii) a first-zero theorem
(Corollary~\ref{cor:firstzero}): for every $f \ge 2$, the first integer zero of
$R(\cdot,f)$ is $k = f$ --- no primality is used --- and (iii) the continuum limit
$R(k,f) \to \sinc^{2}(k/f)$ (Theorem~\ref{thm:continuum}). We close
(Proposition~\ref{prop:firstg}) with the synchronization with the arithmetic
\emph{First-G event}: for $f = \spf(b)$ the smallest prime factor of an integer $b > 1$,
the first integer zero of $R(\cdot,f)$ coincides with the smallest $k$ at which
$\{1,\dots,k\}$ contains a non-coprime element of $\Z/b\Z$. The corrected special value
$\sinc^{2}(1/10) = 25(\sqrt{5}-1)^{2}/(4\pi^{2}) \approx 0.9675312093$ is recorded in
\S\ref{subsec:continuum}; an earlier draft printed the decimal as $0.9355$, a transcription
error that did not affect the closed form. Verification is exhaustive on
$f \in \{3,4,5,6,7,8,9,10,11,12,13,17,19,23\}$ at machine precision (max deviation
$3.33 \times 10^{-16}$); the bundled script \texttt{verify\_J42\_sinc2.py} also confirms
the first-zero theorem for $f \in \{2,\dots,15\}$ and the synchronization for
$b \in \{6,10,15,21,22,35,105\}$.
\end{abstract}

\maketitle

\section*{Lens and substrate}
\label{sec:lens}

This paper works on $\Z/N\Z$ as the cyclic-group quantum-mechanics substrate, with no
TSML/BHML lens dependence. The momentum--position duality is the standard one (DFT-related
orthonormal bases). The theorems are statements about the Fej\'er-kernel restriction to
integer arguments and its arithmetic synchronization with smallest-prime-factor
partitions. No 10-operator framework is invoked; the substrate $\Z/10\Z$ enters only as a
special case ($f = 10$ in \S\ref{subsec:continuum}) and not load-bearingly. The note sits
adjacent to the TIG family-of-finite-magmas program \cite{SandersGishFourCore} and
borders, but does not overlap, the corpus on small finite commutative non-associative
structures pioneered by Dr\'apal and Wanless \cite{DrapalWanless2021}; the present
identity is associative (Fej\'er on $\Z/N\Z$) and stands on its own.

\tableofcontents

%==============================================================
\section{Introduction}
\label{sec:intro}
%==============================================================

Position--momentum duality on the cyclic group $\Z/N\Z$ is the natural arena of finite QM:
the position basis $\{|j\rangle : j \in \Z/N\Z\}$ and the momentum basis
$\{|\hat p\rangle : \hat p \in \Z/N\Z\}$ are related by the discrete Fourier transform
\[
  \langle j \mid \hat p\rangle \;=\; \tfrac{1}{\sqrt{N}}\, e^{2\pi i j \hat p / N}.
\]
A natural physical question is: if we project a momentum eigenstate $|\hat p\rangle$ onto
a position-space rectangular window of size $k$ (the indicator of $\{1,2,\dots,k\}$), what
is the squared overlap with the \emph{normalized} window state?

The answer is the discrete $\sinc^{2}$ function $R(k,f)$ studied here, where
$f = N/\hat p$ is the momentum-quantum index. We present a clean, self-contained
derivation, the elementary closed form (a one-line geometric-sum argument), and the key
consequences for finite QM. The arithmetic counterpart --- the synchronization with the
First-G event in coprimality partitions, \S\ref{sec:firstg} --- makes this object also a
useful reference for number-theoretically flavored finite QM.

\paragraph{Plan.} Section~\ref{sec:setup} sets up the finite-QM Hilbert space and defines
$R(k,f)$. Section~\ref{sec:closedform} proves the closed form
(Theorem~\ref{thm:closedform}). Section~\ref{sec:qm} derives the QM-relevant
consequences: squared overlap with the normalized window (Proposition~\ref{prop:overlap}),
first-zero theorem (Corollary~\ref{cor:firstzero}), and continuum limit
(Theorem~\ref{thm:continuum}). Section~\ref{sec:firstg} records the synchronization with
the First-G event. Section~\ref{sec:tier} collates status, lens scope, and per-venue cap
notes.

%==============================================================
\section{Finite QM on $\Z/N\Z$}
\label{sec:setup}
%==============================================================

\paragraph{Position and momentum bases.}
The cyclic group $\Z/N\Z$ supports a Hilbert space $\mathcal{H} = \C^{N}$ with orthonormal
position basis $\{|j\rangle\}_{j=0}^{N-1}$ and orthonormal momentum basis
$\{|\hat p\rangle\}_{\hat p=0}^{N-1}$ related by
\[
  |\hat p\rangle \;=\; \tfrac{1}{\sqrt{N}}\sum_{j=0}^{N-1} e^{2\pi i j \hat p / N}\,
  |j\rangle.
\]

\paragraph{Position-window state.}
For $1 \le k \le N$, the rectangular position window of size $k$ is the unnormalized
state
\[
  |W_{k}\rangle \;=\; \sum_{j=1}^{k} |j\rangle, \qquad \langle W_{k} \mid W_{k}\rangle = k.
\]
The normalized window is $|w_{k}\rangle = k^{-1/2}\,|W_{k}\rangle$.

\paragraph{Squared overlap.}
The squared overlap of a momentum eigenstate $|\hat p\rangle$ with the normalized window
$|w_{k}\rangle$ is
\[
  |\langle \hat p \mid w_{k}\rangle|^{2}
  \;=\; \tfrac{1}{N k}\Bigl|\sum_{j=1}^{k} e^{-2\pi i j \hat p / N}\Bigr|^{2}.
\]
For the canonical case $\hat p / N = 1/f$ (equivalently $\hat p = N/f$, requiring $f \mid
N$), this is
\begin{equation}\label{eq:overlap-normalized}
  |\langle \hat p \mid w_{k}\rangle|^{2}
  \;=\; \tfrac{k}{N}\, R(k,f), \qquad
  R(k,f) \;=\; \tfrac{1}{k^{2}}\Bigl|\sum_{j=1}^{k} e^{2\pi i j / f}\Bigr|^{2}.
\end{equation}
We focus on $R(k,f)$ as the dimensionless kernel, with the QM normalization factor $k/N$
separated.

%==============================================================
\section{The closed form}
\label{sec:closedform}
%==============================================================

\begin{theorem}[Discrete $\sinc^{2}$ identity]
\label{thm:closedform}
For every integer $f \ge 2$ and every integer $k \ge 1$,
\[
  R(k,f) \;=\; \frac{\sin^{2}(\pi k/f)}{k^{2}\sin^{2}(\pi/f)}.
\]
\end{theorem}

\begin{proof}
The geometric series gives
\[
  \sum_{j=1}^{k} e^{2\pi i j/f}
  \;=\; e^{2\pi i / f}\,\frac{1 - e^{2\pi i k/f}}{1 - e^{2\pi i/f}}.
\]
Taking the squared modulus and dividing by $k^{2}$, with
$|1 - e^{i\theta}|^{2} = 4\sin^{2}(\theta/2)$,
\[
  R(k,f) \;=\; \frac{1}{k^{2}} \cdot \frac{4\sin^{2}(\pi k/f)}{4\sin^{2}(\pi/f)}
  \;=\; \frac{\sin^{2}(\pi k/f)}{k^{2}\sin^{2}(\pi/f)}. \qedhere
\]
\end{proof}

\begin{remark}[Fej\'er-kernel attribution]
The closed form is the Fej\'er kernel $K_{n}(\theta)$ at $\theta = 2\pi/f$, $n = k$,
normalized so that $R$ has unit value at $k = 0$ in the limit and matches $\sinc^{2}$ in
the continuum (\S\ref{subsec:continuum}). The identity dates to Fej\'er's 1900 paper
\cite{Fejer1900}; the contribution of this note is the QM-on-cyclic-group reading
(\S\S\ref{sec:setup},~\ref{sec:qm}) and the arithmetic-bridge synchronization
(\S\ref{sec:firstg}).
\end{remark}

\begin{remark}[Numerical verification]
The bundled script \texttt{verify\_J42\_sinc2.py} compares the closed form against direct
geometric-sum evaluation on $f \in \{3,4,5,6,7,8,9,10,11,12,13,17,19,23\}$ and
$k \in \{1,\dots,f+1\}$; the maximum deviation is $3.33 \times 10^{-16}$. The grid covers
prime and composite $f$, which is the reason the first-zero theorem of
\S\ref{sec:qm} requires no primality hypothesis.
\end{remark}

%==============================================================
\section{Quantum-mechanical consequences}
\label{sec:qm}
%==============================================================

\subsection{Squared overlap with the normalized position window}

\begin{proposition}[Squared overlap with normalized window]
\label{prop:overlap}
For a momentum eigenstate $|\hat p\rangle$ on $\Z/N\Z$ with $\hat p = N/f$ ($f \mid N$),
the squared transition amplitude with the normalized rectangular position window
$|w_{k}\rangle = k^{-1/2}\sum_{j=1}^{k}|j\rangle$ is
\[
  |\langle \hat p \mid w_{k}\rangle|^{2}
  \;=\; \tfrac{k}{N}\,R(k,f)
  \;=\; \frac{1}{N} \cdot \frac{\sin^{2}(\pi k/f)}{k\,\sin^{2}(\pi/f)}.
\]
\end{proposition}

\begin{proof}
Immediate from \eqref{eq:overlap-normalized} and Theorem~\ref{thm:closedform}.
\end{proof}

\paragraph{Note on interpretation.}
This squared overlap is the \emph{fidelity} between $|\hat p\rangle$ and the normalized
window state $|w_{k}\rangle$, not the probability of position-measurement outcome in
$\{1,\dots,k\}$. The latter (the marginal of the position distribution given the momentum
eigenstate) is the uniform $\sum_{j=1}^{k}|\langle j \mid \hat p\rangle|^{2} = k/N$.
The squared overlap $|\langle \hat p \mid w_{k}\rangle|^{2}$ is a different object --- it
measures coherent alignment between $|\hat p\rangle$ and the unit-norm window state ---
and is the kernel of practical interest for finite-QM scattering and projection problems
on cyclic groups. The un-normalization factor $k$ relating $|W_{k}\rangle$ and
$|w_{k}\rangle$ is explicit, which fixes a misreading in an earlier draft that conflated
this overlap with the position-marginal probability.

\subsection{First-zero theorem}

\begin{corollary}[First zero of $R(\cdot,f)$]
\label{cor:firstzero}
For every $f \ge 2$ and every $k \in \{1,\dots,f-1\}$, $R(k,f) > 0$. At $k = f$,
$R(f,f) = 0$. Hence the first integer zero of $R(\cdot,f)$ in $\Z_{\ge 1}$ is exactly
$k = f$, regardless of whether $f$ is prime or composite.
\end{corollary}

\begin{proof}
For $1 \le k \le f-1$, $f \nmid k$ (since $0 < k < f$), so
$\pi k/f \notin \pi\Z$ and $\sin(\pi k/f) \neq 0$, giving $R(k,f) > 0$ from
Theorem~\ref{thm:closedform}. At $k = f$, $\sin(\pi f/f) = 0$, and the closed form
($0/0$ numerically) is resolved by the geometric sum
$\sum_{j=1}^{f} e^{2\pi i j /f} = 0$, giving $R(f,f) = 0$. The argument does not use the
prime structure of $f$.
\end{proof}

\begin{remark}[Where prime structure enters]
Primality of $f$ is \emph{not} used in Corollary~\ref{cor:firstzero}; the first-zero
result holds uniformly for $f \ge 2$. The prime restriction enters only at the
synchronization with the First-G event (Proposition~\ref{prop:firstg}), where one needs
$\gcd(k,b) = 1$ for $1 \le k < p_{1}$ to deduce that no obstruction lies below the
smallest prime factor.
\end{remark}

\paragraph{QM reading.} The momentum eigenstate $|\hat p\rangle$ with $\hat p = N/f$ has
zero overlap with the normalized window $|w_{f}\rangle$ of size $k = f$. This is the
cyclic-group analog of the first zero of $\sinc^{2}(t)$ at $t = 1$.

\subsection{Continuum limit}
\label{subsec:continuum}

\begin{theorem}[$\sinc^{2}$ continuum limit]
\label{thm:continuum}
Let $\sinc(t) = \sin(\pi t)/(\pi t)$ for $t \neq 0$ and $\sinc(0) = 1$. For every $t > 0$,
\[
  \lim_{\substack{f \to \infty \\ k/f \to t}} R(k,f) \;=\; \sinc^{2}(t).
\]
\end{theorem}

\begin{proof}
From Theorem~\ref{thm:closedform}, $k\sin(\pi/f) = (\pi/f)\cdot k\cdot(1+O(1/f^{2}))
\to \pi t$, so $k^{2}\sin^{2}(\pi/f) \to \pi^{2}t^{2}$. The numerator
$\sin^{2}(\pi k/f) \to \sin^{2}(\pi t)$. Hence
$R(k,f) \to \sin^{2}(\pi t)/(\pi^{2}t^{2}) = \sinc^{2}(t)$.
\end{proof}

\paragraph{Specific exact values (continuum limit at rational arguments).}
The continuum identity $\sinc^{2}(1/2) = (2/3)\cdot 1/\zeta(2) = 4/\pi^{2}
\approx 0.4053$ is the special case $t = 1/2$, where $\sin(\pi/2) = 1$ and the
$\zeta(2) = \pi^{2}/6$ rewriting is the bridge to the Riemann zeta values. At
$t = 1/10$, Ptolemy's pentagon identity $\sin(\pi/10) = (\sqrt{5}-1)/4$ gives the exact
closed form
\[
  \sinc^{2}(1/10) \;=\; \frac{\sin^{2}(\pi/10)}{(\pi/10)^{2}}
  \;=\; \frac{25(\sqrt{5}-1)^{2}}{4\pi^{2}}
  \;\approx\; 0.9675312093.
\]
This is the corrected value. An earlier draft printed the decimal as $0.9355$, a
transcription error; the closed form $25(\sqrt{5}-1)^{2}/(4\pi^{2})$ was always correct,
and the bundled \texttt{verify\_J42\_sinc2.py} confirms agreement between sympy and direct
numerical evaluation to ten digits.

\paragraph{QM reading.}
The continuum limit recovers the rectangular-window momentum spectrum exactly:
$|\langle \hat p \mid \text{window}_{t}\rangle|^{2} \to \sinc^{2}(t)/N$ as the lattice is
refined (with the standard $k/N$ scaling reabsorbed). The discrete $\sinc^{2}$ identity
of Theorem~\ref{thm:closedform} is therefore the natural finite-QM counterpart of the
continuous rectangular-pulse spectrum.

%==============================================================
\section{Synchronization with the First-G event}
\label{sec:firstg}
%==============================================================

The identity of \S\ref{sec:closedform} has a striking arithmetic counterpart. For an
integer $b > 1$, the \emph{coprimality partition} of the alphabet $\{1,\dots,k\}$
relative to $b$ separates units $C_{k}(b) = \{x : \gcd(x,b) = 1\}$ from obstructions
$G_{k}(b) = \{1,\dots,k\}\setminus C_{k}(b)$. The \emph{First-G event} $k^{\star}(b)$ is
the smallest $k$ with $|G_{k}(b)| > 0$.

\begin{proposition}[Synchronization with First-G]
\label{prop:firstg}
For every $b > 1$ with smallest prime factor $p_{1} = \spf(b)$, the First-G event
$k^{\star}(b)$ and the first integer zero of $R(\cdot,p_{1})$ in $\{1,\dots,p_{1}\}$
coincide:
\[
  k^{\star}(b) \;=\; \min\{k \in \{1,\dots,p_{1}\} : R(k,p_{1}) = 0\} \;=\; p_{1}.
\]
\end{proposition}

\begin{proof}
For any $x$ with $1 \le x \le p_{1} - 1$, $\gcd(x,b) = 1$, since the smallest prime
factor of $b$ is $p_{1}$ and no prime $\le p_{1}-1$ divides $b$. Hence
$k^{\star}(b) = p_{1}$. Combined with Corollary~\ref{cor:firstzero} applied at
$f = p_{1}$, this gives the displayed identity.
\end{proof}

\paragraph{Reading.}
The QM-momentum quantum that produces the first position-window zero at $k = p_{1}$ is
precisely the inverse of the smallest prime factor of $b$. This synchronization is the
bridge between cyclic-group QM and the arithmetic structure of $\Z/b\Z$, and is why
$R(k,f)$ appears naturally in finite-dimensional quantum number theory. The arithmetic
side --- including the CRT-based inheritance of the zero structure through squarefree
moduli --- is treated in the companion paper \cite{SandersGishSincZero} (J04); the
present note restates the synchronization in the QM context.

%==============================================================
\section{Tier, status, lens scope, per-venue cap}
\label{sec:tier}
%==============================================================

\subsection{PROVED / COMPUTED / STRUCTURAL RHYME / OPEN}

\begin{itemize}\setlength\itemsep{2pt}
  \item \textbf{PROVED.} Theorem~\ref{thm:closedform} (closed form, Fej\'er kernel at
    $\theta = 2\pi/f$, $n = k$); Proposition~\ref{prop:overlap} (squared overlap with
    normalized window); Corollary~\ref{cor:firstzero} (first zero at $k = f$ for every
    $f \ge 2$, no primality used); Theorem~\ref{thm:continuum} ($\sinc^{2}$ continuum
    limit); Proposition~\ref{prop:firstg} (synchronization
    $k^{\star}(b) = \spf(b)$ with the First-G event).
  \item \textbf{COMPUTED.} \texttt{verify\_J42\_sinc2.py}: max deviation
    $3.33\times 10^{-16}$ for Theorem~\ref{thm:closedform} over
    $f \in \{3,4,5,6,7,8,9,10,11,12,13,17,19,23\}$ and $k \in \{1,\dots,f+1\}$;
    unrestricted Corollary~\ref{cor:firstzero} confirmed for $f \in \{2,\dots,15\}$;
    corrected $\sinc^{2}(1/10) = 25(\sqrt{5}-1)^{2}/(4\pi^{2}) \approx 0.9675312093$
    (sympy and numpy agreement to ten digits); Proposition~\ref{prop:firstg} confirmed
    for $b \in \{6, 10, 15, 21, 22, 35, 105\}$.
  \item \textbf{STRUCTURAL RHYME.} The closed form $25(\sqrt{5}-1)^{2}/(4\pi^{2})$ at
    $f = 10$ folds in Ptolemy's pentagon identity
    $\sin(\pi/10) = (\sqrt{5}-1)/4$, connecting cyclic-group QM at $f = 10$ to
    regular-pentagon arithmetic. Cited as a structural motif, not load-bearing.
  \item \textbf{OPEN.} Does $R(k,f)$ for $f$ ranging over the prime factors of squarefree
    $b$ encode the full coprimality structure of $(\Z/b\Z)^{\times}$? The CRT-based
    inheritance result of \cite{SandersGishSincZero} treats the squarefree case; the
    general case is open.
\end{itemize}

\subsection{Lens scope.}
This paper carries no TSML / BHML lens dependence. The mathematical content is
finite-dimensional Fourier analysis on the cyclic group, with an explicit QM
interpretation. The lens-ownership paragraph is at the top (\S\ref{sec:lens}).

\subsection{Per-venue cap.}
This is the third \emph{Journal of Mathematical Physics} target in the J-series after
J40 (BB Bridge) and J41 (YM Mass Gap Bridge); the 2/quarter cap is reached. We submit
to \emph{Letters in Mathematical Physics} as the preferred venue; \emph{Journal of
Physics A: Mathematical and Theoretical} and \emph{Communications in Mathematical
Physics} are documented alternatives in the cover letter. See cover letter and README
\S5 for details.

\subsection{Tier.}
Central claims are Tier 1/2 (clean theorems, fully proved). The closed form is
elementary; the QM consequences are direct corollaries; the synchronization is one
combination step from companion work \cite{SandersGishSincZero}.

%==============================================================
\section*{Acknowledgments}
%==============================================================

This note extracts in self-contained form the QM-on-cyclic-group reading of the discrete
$\sinc^{2}$ identity descending from Theorem~3.1 of the companion First-G / Zero-Law
manuscript \cite{SandersGishSincZero}. We thank the colleagues who pointed out the
$\sinc^{2}(1/10)$ decimal transcription error in an earlier draft (the closed form was
always correct; only the printed decimal was wrong; see
\S\ref{subsec:continuum}). Algorithmic verification scripts were developed with
assistance from Anthropic Claude sessions in 2026; all theorems and numerical checks were
re-verified independently and are deposited at
\url{https://github.com/TiredofSleep/ck/tree/tig-synthesis}.

%==============================================================
\begin{thebibliography}{99}
%==============================================================

\bibitem{Fejer1900}
L.~Fej\'er,
\emph{Sur les fonctions born\'ees et int\'egrables},
\textit{Math.\ Ann.} \textbf{58} (1900), 51--69.

\bibitem{Zygmund2002}
A.~Zygmund,
\emph{Trigonometric Series}, 3rd ed.,
Cambridge University Press, 2002.

\bibitem{Schwinger1960}
J.~Schwinger,
\emph{Unitary operator bases},
\textit{Proc.\ Nat.\ Acad.\ Sci.} \textbf{46} (1960), 570--579 and 893--897.

\bibitem{Wootters1987}
W.~K.~Wootters,
\emph{A Wigner-function formulation of finite-state quantum mechanics},
\textit{Ann.\ Phys.} \textbf{176} (1987), 1--21.

\bibitem{Vourdas2004}
A.~Vourdas,
\emph{Quantum systems with finite Hilbert space},
\textit{Rep.\ Prog.\ Phys.} \textbf{67} (2004), 267--320.

\bibitem{Vourdas2017}
A.~Vourdas,
\emph{Finite and Profinite Quantum Systems},
Springer, 2017.

\bibitem{Apostol1976}
T.~M.~Apostol,
\emph{Introduction to Analytic Number Theory},
Springer, 1976.

\bibitem{HardyWright2008}
G.~H.~Hardy, E.~M.~Wright,
\emph{An Introduction to the Theory of Numbers}, 6th ed.,
Oxford University Press, 2008.

\bibitem{DrapalWanless2021}
A.~Dr\'apal, I.~M.~Wanless,
\emph{On the number of quasigroups and Latin squares},
\textit{J.\ Combin.\ Theory Ser.\ A} \textbf{184} (2021), 105510.
(Companion in the small-finite-commutative-non-associative neighborhood --- opposite
extremum: maximally non-associative.)

\bibitem{SandersGishSincZero}
B.~R.~Sanders, M.~Gish,
\emph{Full-Period Cancellation of $R(k,f)$ and the spf-Localization for Squarefree
Moduli} (J04),
submitted to \textit{Integers}, 2026.

\bibitem{SandersGishFourCore}
B.~R.~Sanders, M.~Gish,
\emph{Joint Closure, Per-Coordinate Fuse Data, and a Closed-Form Algebraic Attractor of
Two Commutative Binary Operations on $\Z/10$},
submitted to \textit{Journal of Algebra}, 2026.

\bibitem{TIGsynthesis2026}
B.~R.~Sanders,
\emph{Trinity Infinity Geometry Synthesis},
\url{https://github.com/TiredofSleep/ck},
DOI: 10.5281/zenodo.18852047, 2026.

\end{thebibliography}

\end{document}
